% !PRINTED 169
\weberpage{169}
%<<CONT
further, the system $\mathfrak{S}$ contains forms of degree 0, that is, constants different from zero, then our theorem is self-evident, since we may take such a constant for $F_1$, and, if $A_1 = F : F_1$ is set, $F = A_1F_1$. In this form, however, the theorem is without content. But if we exclude from $\mathfrak{S}$ all constant forms, there remains a system $\mathfrak{S}'$ which no longer contains any forms of degree 0, and for which our theorem likewise holds. The basis $F_1, F_2, \ldots F_\mu$ then likewise contains no forms of degree 0, and we can therefore also express Theorem I. thus:

\gesperrt{II. All forms of the system $\mathfrak{S}$ of positive degree can be expressed by means of a basis $F_1, F_2, \ldots F_\mu$, whose elements are of positive degree, in the form:}
\weq{1}{F = \Phi_1F_1 + \Phi_2F_2 + \cdots + \Phi_\mu F_\mu.}
\gesperrt{The degrees of $\Phi_1, \Phi_2, \ldots \Phi_\mu$ are then lower than the degree of $F$.}

This theorem finds its most important applications in investigations of the possibility of representing all forms of a certain system $\mathfrak{S}$ as \gesperrt{integral rational functions of a finite number} among them. Such a system of forms is called a \gesperrt{finite} one (not in the sense that it consists of only a finite number of forms). The following theorem holds:

\gesperrt{III. If the coefficients $\Phi_1, \Phi_2, \ldots \Phi_\mu$ in the representation (1) of Theorem II. can be so chosen for every form $F$ in $\mathfrak{S}$ that, when they are not constant, they themselves belong to the system $\mathfrak{S}$, then the system $\mathfrak{S}$ is finite.}

This follows immediately from the fact that the degrees of the forms $\Phi_1, \Phi_2, \ldots \Phi_\mu$ are lower than the degree of $F$. If, then, one applies the representation (1) to the non-constant ones among the functions $\Phi$, one arrives at coefficients of still lower degree, and so, on repeated application of this procedure, one must finally come to constants.

From this there now follows, by a very simple mode of reasoning, which I owe to an oral communication from \gesperrt{Hurwitz}, the finiteness of the system of invariants $\mathfrak{J}$ of a finite group of linear substitutions $S$.
